\documentclass[../main.tex]{subfiles}
\begin{document}
%==========================================TIÊU ĐỀ TẬP========================================
\begin{table}[h]
    \begin{adjustwidth}{-1cm}{}
        \begin{tabular}{>{\centering\arraybackslash}p{6cm}|>{\raggedright\arraybackslash}p{12.5cm}}
            \multirow{5}{6cm}
            % Cột bên trái: ÉPISODE và số tập
            {\\[-40pt]
                \flaregothic\fontsize{20pt}{20pt}\selectfont \centering
                \color{eptype}SPÉCIAL\color{black}\\
                \burbank\fontsize{72pt}{72pt}\selectfont
                \color{epnum}IX\color{black}\\[6pt]
                \cafeteria\fontsize{20pt}{20pt}\selectfont\color{epnum}(S-9)\color{black}
                \\[18pt]
                \sen\fontsize{18pt}{18pt}\selectfont
                \color{epattr}Mini-histoire
                \color{black}
            }
             &
            % Dòng thứ nhất: tên tập tiếng Nhật
            {\fotskip\fontsize{18pt}{18pt}\selectfont ア マ ネ カ ナ ッ タ
            }                               \\[6pt]
            % Dòng thứ hai: tên tập tiếng Anh
             & {\cafeteria\fontsize{18pt}{18pt}\selectfont A Gossip about Amane Kanata}          \\[4pt]
            % Dòng thứ ba: tên tập tiếng Việt
             & {\vnmsans\fontsize{18pt}{18pt}\selectfont Những câu chuyện phiếm về Amane Kanata} \\[8pt]
            % Dòng thứ tư: Nhân vật xuất hiện
             & {\caviar{\fontsize{14pt}{14pt}\selectfont Nhân vật: Theo từng câu chuyện}}
            \\[8pt]
            % Dòng cuối cùng: Ngày phát hành
             & {\continuum\fontsize{16pt}{16pt}\selectfont Ngày phát hành: 10/07/2020}           \\[18pt]
        \end{tabular}
    \end{adjustwidth}
\end{table}
%=========================================TRUYỆN CHÍNH========================================
\color{black}

% Hộp tóm tắt
\txtbx[2]{{{\color{vol} \uline{Tóm tắt:}} Bảy câu chuyện nhỏ về cuộc sống của nàng thiên sứ Amane Kanata tại \hll. Từ chiếc áo phông bị gọi là "giẻ lau", đến việc cô bỗng dưng biến thành vẹt, rồi vầng hào quang bỗng dưng to như phi tiêu và quay điên cuồng, và cả trong ván Ma Sói nơi cô bị cả hội chỉ điểm là Ma Sói dù đang cố giấu, và cuối cùng bị loại vì ``linh cảm của loài Chó''. Cô còn tự tổ chức cuộc bình chọn nhân vật yêu thích với năm vị trí đều là chính mình. Phần trailer hoành tráng cho tập sau với chủ đề ``CÁI CHẾT CỦA KANATA'' kết thúc bằng bữa ăn ramen ấm áp, nơi Coco và mọi người cổ vũ Kanata trước buổi hologra sắp tới.}}

\begin{center}
    \uline{(Trước khi Kanata thực sự đến văn phòng và gặp mặt Hikawa Kuon, cô nàng cũng đã có một buổi phát trực tiếp để ra mắt mô hình 3D đầu tiên của mình. Trong buổi phát sóng đó, cô nàng còn có một thước phim ngắn để chào mừng sự kiện này.\\Tập đặc biệt này sẽ lấy trực tiếp từ thước phim đó, pha trộn các thành viên từ hai phiên bản: giọng gốc và giọng thay thế.)}
\end{center}

\img{s-8a}

\pagebreak

\subsection*{\phantomsection\label{ep-S7.I}}
\addcontentsline{toc}{subsection}{{\sen{-Histoire 1-}} {\caviar Tấm giẻ lau}}
\stpt{-Histoire 1-}
\stcap{Tấm giẻ lau}
\begin{center}
    {\caviar\fontsize{12pt}{12pt}\selectfont Nhân vật: \emp{kuon1} \emp{haato} \emp{coco} \emp{kanata}}
\end{center}

Không khí trong văn phòng hôm nay thật sự yên bình. Một cái yên bình hiếm hoi khiến người ta có cảm giác rằng, có lẽ, chỉ có lẽ thôi, \hll\ cũng có thể có những ngày không có chuyện gì xảy ra. Coco đang cắm đầu vào máy tính với đống báo cáo còn tồn đọng, Haato đang lướt điện thoại với vẻ mặt thư thái, còn tôi thì đang cố gắng sắp xếp lại đống giấy tờ mà không muốn mất hết hy vọng vào cuộc sống.

Rồi cánh cửa bật tung.

Không phải mở, mà là bật tung. Một lực mạnh bất thường khiến nó va vào tường với một tiếng *RẦM* lớn. Và từ sau cánh cửa, Kanata lao vào như một cơn gió nhiệt đới, vầng hào quang trên đầu xoay tít, đôi cánh nhỏ phe phẩy, mắt sáng rực như thể vừa phát minh ra thứ gì đó vĩ đại: ``Mọi người ơi, nhìn này, nhìn này! Em có thứ này!''

Ba cái đầu trong phòng đồng loạt ngẩng lên. Coco ngừng gõ bàn phím, Haato ngừng lướt điện thoại, tôi ngừng ảo tưởng về một ngày yên bình.

Coco cất giọng hỏi, vừa tò mò vừa có phần dè dặt, vì kinh nghiệm cho thấy mỗi khi Kanata phấn khích, thường có chuyện kỳ quặc xảy ra: ``Gì thế, Kanatan? Lại tự làm thêm món đồ chơi mới à?''

Tôi cũng theo đó lên tiếng thắc mắc, đặt bút xuống, ngả người ra ghế: ``Có cái gì mà em phấn khích thế? Mới sáng ra đã thấy em nhảy nhót như vậy, chắc là có tin vui?''

Kanata hớn hở, hai tay giơ lên trời như đang chuẩn bị cho một phép màu: ``Em đã tự tay làm một sản phẩm mới! Lần này em đầu tư công sức lắm đấy!''

Vừa nói, cô vừa giơ cao một chiếc áo phông trắng. Trên đó có dòng chữ to tướng ``PP-Tenshi'' được viết tay bằng mực vải. Nét chữ còn hơi nguệch ngoạc, màu sắc cũng không đều lắm, nhưng có vẻ như cô đã rất tâm huyết với sản phẩm này.

\img{s-8b}

``Đây! Em đã tự thiết kế nó! Trông nó dễ thương chứ? Em đã mất cả buổi tối hôm qua để vẽ, đến nỗi tay cứng đơ luôn đấy!''

Haato ngắm nghía chiếc áo một hồi, rồi thốt lên với giọng đầy thiện chí: ``Trông nó đáng yêu thật đấy nhỉ! Chữ viết độc đáo quá, giống kiểu nghệ thuật đường phố ấy!''

Tôi cũng gật gù, định an ủi cô bé thiên sứ đang tỏa ra một vầng hào quang của sự tự hào: ``Anh thấy sáng tạo đấy. Anh đánh giá cao về chuyện này. Ít nhất thì em cũng không đổ bột mì lên người như lần trước.''

Nghe lời khen từ hai người, Kanata càng trở nên phấn chấn hơn. Cô quay sang Coco, người vẫn đang im lặng từ nãy đến giờ với vẻ mặt đầy hy vọng: ``Đúng chứ!? Coco, bà thấy sao? Bà có muốn một cái không? Tôi có thể làm cho bà một cái với màu sắc khác!''

Không giống như phản ứng tích cực của Haato và tôi, Coco chỉ lặng thinh nhìn chiếc áo với một ánh mắt khó diễn tả. Nó trông chẳng phải là khinh bỉ, cũng chẳng phải thích thú, mà đó là một cái nhìn giống như khi một người bình thường nhìn thấy một món đồ ở cửa hàng và tự hỏi ``ai lại mua cái này nhỉ?''. Sau vài giây im lặng đến khó chịu, nàng Rồng nhướng mày, hờ hững phán một câu khiến cả ba người còn lại chết sững: ``Cái giẻ lau bẩn thỉu này là sao đây? Trông như ai đó đã giặt quần áo xong rồi dùng nó để lau sàn vậy.''

Tôi và Haato đồng thanh, mắt mở to: ``Hả? Giẻ lau? Ý em là sao?''

Kanata - tác giả đang đứng trước thành quả của mình - bất chợt nổi cáu, trái ngược hoàn toàn với sự mong chờ ban nãy: ``\textbf{Không phải!} Bà bảo cái áo tôi làm là giẻ lau á? Tôi đã mất bao công sức để vẽ nó đấy!''

Haato đổ mồ hôi hột, vội vã đứng ra làm trung gian: ``Coco-chan, đừng có thô lỗ thế chứ-chama! Dù gì Kanata-chan cũng đã làm ra nó rồi mà! Nói nó là giẻ lau thì hơi quá đấy!''

Nhưng Coco vẫn tỉnh bơ như không, nhún vai: ``Thì người ta nói vậy nên tôi nói lại thôi. Sao phải căng thế? Mà nhìn kỹ thì nó cũng có công dụng lau bàn đấy chứ.''

Tôi nhìn nàng Rồng với ánh mắt bất lực, định lên tiếng giải thích, nhưng chưa kịp mở miệng thì Kanata đã nổi đóa. Cô lập tức lao đến bàn mình, giật lấy chiếc máy tính xách tay, mở tung lên với một tiếng ``rầm'': ``Được rồi! Tôi sẽ cho bà xem thế nào mới thực sự là giẻ lau! Bà sẽ thấy trên internet người ta nói gì về nó!''

Cô hậm hực mở trình duyệt lên, gõ vào thanh tìm kiếm của PP Search dòng chữ ``ボロ雑巾'' (giẻ rách). Ngón tay cô bấm phím Enter với đầy quyết tâm, như thể đang ký một bản án.

Và kết quả trả về trên màn hình chính là hình ảnh chiếc áo phông của cô, đúng cái áo mà cô vừa giơ lên cho ba người xem. Cùng một dòng chữ, cùng một kiểu chữ, cùng một màu mực. Nó được đăng tải trên một trang web với dòng mô tả: ``Áo thun giẻ lau - độc đáo, thời trang, và\dots\ lau sạch.''

\img{s-8c}

Cả phòng im bặt.

Mắt Kanata mở to hết cỡ, mặt cô đờ ra như thể thế giới quan vừa bị sụp đổ. Miệng cô há hốc, nhưng không phát ra âm thanh nào. Haato che miệng cố nhịn cười, vai run lên bần bật. Coco thì chỉ khoanh tay gật gù như thể đã đoán trước kết quả này từ đầu.

Tôi thở dài, đưa tay xoa trán, cố tìm cách vớt vát lại tình hình: ``Anh nghi chắc mọi người gọi như thế chỉ là đùa thôi, chứ anh không nghĩ họ có ý gì xấu đâu\dots\ Có thể họ thấy nó độc đáo nên mới so sánh vậy.''

Nhưng Kanata đã chẳng còn nghe thấy gì nữa. Cô run rẩy nắm chặt nắm đấm, vẻ mặt đầy phẫn uất như thể muốn đập nguyên cái máy tính xuống đất, rồi phanh phui cả internet. Haato và tôi phải vội vàng giữ chặt cô lại, một người giữ tay, một người giữ vai.

Trong khi đó, Coco chỉ bình tĩnh nhấp một ngụm nước, đặt cốc xuống, thản nhiên kết luận như thể vừa khẳng định một chân lý hiển nhiên: ``Ừ, đúng là đồ giẻ lau thật mà. Internet không bao giờ sai.''

Và thế là, văn phòng chìm trong hỗn loạn khi nàng Thiên sứ rượt đuổi nàng Rồng khắp nơi, còn Haato và tôi thì vất vả ngăn cản một vụ đập phá đồ đạc lớn nhất trong tháng.

Tiếng bước chân, tiếng la hét, tiếng cười, và một chiếc áo phông bị vứt lại trên bàn, im lìm như một nhân chứng cho sự sụp đổ của một ước mơ. Tôi nhìn chiếc áo, rồi nhìn Coco đang chạy vòng quanh bàn để tránh Kanata, và thở dài: ``\textit{Đúng là một ngày bình yên\dots}''

\subsection*{\phantomsection\label{ep-S7.II}}
\addcontentsline{toc}{subsection}{{\sen{-Histoire 2-}} {\caviar Vẹt}}
\stpt{-Histoire 2-}
\stcap{Vẹt}
\vspace{-4pt}
\begin{center}
    {\caviar\fontsize{12pt}{12pt}\selectfont Nhân vật: \emp{kuon1} \emp{haato} \emp{coco} \emp{kanata}}
\end{center}

Một buổi chiều yên ả hiếm hoi khác, bốn người ngồi quây quần trong văn phòng. Không có bom, không có đột kích, không có khủng long. Chỉ là một khoảng lặng quý giá giữa những cơn bão.

Coco ngồi gác chân lên bàn, tay cầm cốc trà sữa, nhìn ra cửa sổ. Haato đang cặm cụi viết gì đó. Tôi thì đang cố gắng tìm cách sắp xếp lại đống giấy tờ mà không muốn bỏ cuộc. Kanata ngồi im, có vẻ đang suy nghĩ về điều gì đó sâu xa.

Rồi cô lên tiếng với một vẻ mặt trầm tư, như một triết gia đang đặt câu hỏi về bản thể: ``Mọi người hay nói rằng tôi hay lặp lại lời người ta nói lắm\dots''

\img{s-8e}

Ngay lập tức, Coco bắn trả không chút do dự, giọng như thể đã chờ đợi câu hỏi này từ lâu: ``Thì đúng quá rồi còn gì nữa! Nãy giờ tôi đang tự hỏi khi nào cô mới nhận ra đấy.''

Haato ngạc nhiên với lời bình của nàng Rồng, giọng đầy thiện chí: ``Thật thế á? Nhưng chị tưởng là hầu như ai cũng đôi khi lặp lại lời mình chứ? Đôi khi mình nói một câu, rồi lại nói lại để nhấn mạnh\dots''

Tôi vội lên tiếng đính chính, cố gắng làm rõ sự khác biệt: ``Cái đó gọi là nói lắp. Nhưng như Haato nói đấy, ai mà chẳng từng bị như vậy một, hai lần. Đặc biệt là khi em đang cố gắng nhớ lại mình đã nói gì.''

``Vấn đề là bà ấy không phải chỉ đôi khi, mà là suốt ngày như thế!'', Coco cũng nói góp, tay vung vẩy như thể đang minh họa một phát hiện khoa học, ``Cô ấy lặp lại mọi thứ! Nếu tôi nói `chào', cô ấy sẽ nói `chào'. Nếu tôi nói `tôi đói', cô ấy sẽ nói `tôi đói'. Nếu tôi nói `tôi là Rồng', cô ấy sẽ nói `tôi là Rồng' - mặc dù cô ấy không phải là Rồng.''

Tới lúc này, Kanata gật đầu thừa nhận, giọng thành thật: ``Ừ\dots\ Tôi toàn bị như thế suốt.''

Tôi nhân cơ hội chỉ điểm ngay khi cô vừa dứt câu: ``Như vừa rồi chẳng hạn. Em vừa lặp lại chính xác suy nghĩ của anh sau khi anh đã nói xong.''

Và Kanata, như một phản xạ tự nhiên, lặp lại đúng ý mà tôi vừa nói: ``À, như vừa rồi\dots''

Ba người còn lại nhìn nhau. Một khoảng lặng ngắn, nặng nề, đầy ý nghĩa. Rồi Coco đột nhiên hỏi một câu khiến cả phòng lặng như tờ, giọng đầy vẻ tò mò của một nhà khoa học phát hiện ra loài mới: ``Bà có phải là con vẹt không vậy?''

Haato và tôi lập tức đồng tình, không hẹn mà gặp:

``Ừ, chị cũng thấy thế. Chẳng khác gì mấy con vẹt ở cửa hàng thú cưng ấy.''

``Anh cũng nghĩ vậy. Em mà có lông, chắc người ta đem em ra bán mất.''

Kanata thoáng giật mình, đôi mắt mở to, rồi cô tự hỏi với vẻ mặt cực kỳ nghiêm túc: ``Tôi có phải là con vẹt không ta?''

Ngay giây phút đó, cả ba người còn lại đồng loạt đổ mồ hôi hột. Một giọt mồ hôi lăn trên trán tôi, và tôi tự hỏi liệu câu hỏi đó có phải là một phép màu hay một lời nguyền. Kanata chớp chớp mắt, đăm chiêu suy nghĩ, vầng hào quang trên đầu xoay chậm, rồi đột nhiên, một luồng ánh sáng lóe lên xung quanh cô.

Nó không chói lòa, nhưng đủ mạnh để tôi phải nheo mắt. Khi ánh sáng tan đi, trước mặt ba người không còn là Kanata nữa.

Thay vào đó là một con vẹt xám. Nó có một chỏm tóc trắng trên đầu, đôi mắt đen láy, và vầng hào quang vẫn lơ lửng phía trên, nhưng giờ nó trông như một chiếc vòng cổ trang trí. Con vẹt này trông y như một phiên bản chim của nàng Thiên sứ.

\img{s-8d}

Ba người nhìn chằm chằm con vẹt trước mặt, không nói nên lời. Coco bỏ cốc trà xuống, mắt mở to. Haato thì đưa tay che miệng. Tôi thì chỉ biết đứng hình, tự hỏi cuộc sống của mình đã đi đến đâu.

Họ cùng nhau đồng thanh: ``Chuẩn con vẹt thật rồi.''

Kanata (hay đúng hơn là con vẹt Kanata) nghiêng đầu một góc 90 độ, chớp mắt nhìn họ, rồi cất giọng, vẫn là giọng của cô nhưng vang lên từ một chiếc mỏ chim: ``Chuẩn con vẹt thật rồi.''

Cả phòng im lặng. Chính cô nàng ấy đã thực sự thừa nhận như vậy, cả về nghĩa đen và nghĩa bóng. Cô vừa lặp lại lời của họ, và cũng vừa xác nhận rằng mình đúng là một con vẹt.

Coco bịt trán, ngước lên trời, giọng đầy bất lực: ``Xong rồi, từ nay văn phòng này có một con vẹt thật sự rồi. Ai mà tin được chứ?''

Haato, với đầu óc kinh doanh của mình, mắt sáng lên như sao, như thể vừa nhìn thấy một cơ hội đầu tư mới: ``Liệu có bán được với giá cao không nhỉ? Vẹt biết nói, lại biến hình từ thiên sứ\dots\ chắc chắn sẽ là mặt hàng hot trên thị trường.''

Ngay lập tức, tôi lên tiếng phản bác ngay, giọng đầy bất lực: ``Đừng có biến đồng nghiệp thành hàng hóa chứ! Dù cô ấy có biến thành vẹt, thì vẫn là thành viên của \hll\.''

Trong khi cả ba người còn đang tranh cãi - Coco đang tính xem sẽ nuôi nó ở đâu, Haato đang tìm giá thị trường, còn tôi thì cố gắng giữ trật tự - vẹt Kanata vỗ cánh bay lên. Nó quạt gió tứ phía, bay vòng quanh phòng, rồi hô lớn bằng giọng trong trẻo, lặp đi lặp lại như một câu thần chú: ``Tôi có phải là con vẹt không ta! Tôi có phải là con vẹt không ta!''

Chẳng ai biết hôm nay cô ấy sẽ biến trở lại hình dạng ban đầu hay không. Có thể cô sẽ thức dậy vào ngày mai và lại là một thiên sứ bình thường, hoặc có thể cô sẽ mãi mãi là một con vẹt xám với vầng hào quang trên đầu.

Nhưng ít nhất, một điều đã chắc chắn: từ nay, biệt danh ``PP-Vẹt'' của cô sẽ đi vào huyền thoại.

\subsection*{\phantomsection\label{ep-S7.III}}
\addcontentsline{toc}{subsection}{{\sen{-Histoire 3-}} {\caviar Vầng quang của Kanata}}
\stpt{-Histoire 3-}
\stcap{Vầng quang của Kanata}
\vspace{-4pt}
\begin{center}
    {\caviar\fontsize{12pt}{12pt}\selectfont Nhân vật: \emp{kuon1} \emp{suisei} \emp{aqua} \emp{kanata}}
\end{center}

Tôi đang sắp xếp lại đống tài liệu, Coco đã về nhà sớm, Haato đang đi mua đồ ăn. Chỉ còn tôi và Kanata, và cô ấy đang say sưa với điện thoại, thỉnh thoảng cười khúc khích một mình.

Thế rồi, bầu không khí bỗng chốc bị phá vỡ bởi một tiếng hét thất thanh: ``\textbf{Cứu tôi với!!}''

Tôi ngẩng đầu lên, và cảm giác bình yên tan biến ngay lập tức. Kanata đang chạy vòng quanh phòng, hai tay vung vẩy như thể vừa nhìn thấy ma. Nhưng không phải ma, mà phía sau nàng Thiên sứ, một vầng hào quang khổng lồ đang quay tít với tốc độ kinh hoàng. Trông nó chẳng khác gì một chiếc phi tiêu máy bay khổng lồ được gắn động cơ phản lực, sẵn sàng cắt bất cứ thứ gì đứng gần.

Tôi vừa né tránh khi nó lao về phía tôi, vừa cuống quýt tìm cách đưa nó về trạng thái bình thường: ``Từ từ, bình tĩnh! Để anh thử xem\dots\ Em có biết cách tắt nó không? Hay có nút điều chỉnh nào không?''

Kanata vừa chạy vừa đáp, giọng như sắp khóc: ``Em không biết nữa! Nó cứ tự nhiên to ra và quay như điên!''

Ngay lúc đó, cánh cửa mở ra. Aqua và Suisei - hai người vừa mới tới - tò mò tiến lại gần, quan sát cảnh tượng hỗn loạn trước mắt.

\img{s-8f}

Nàng Hầu gái cất giọng hỏi, vừa tò mò vừa có phần thích thú: ``Kanata-chan đang làm cái gì vậy? Tập thể dục à? Hay đang tập một tiết mục mới cho buổi hòa nhạc?''

Nàng Sao chổi phấn khích hơn, mắt sao chổi của cô sáng rực như thể vừa nhìn thấy một trò chơi mới: ``Cái ngôi sao đang quay sau lưng là thế nào đây? Trông có vẻ vui nha! Nó có thể dùng để cắt bánh mì không?''

Nghe vậy, mặt tôi đanh lại, giọng đầy bất lực: ``Nhìn chả vui vẻ gì đâu\dots\ Nó vừa cắt mất một góc bàn làm việc của anh đấy.''

Kanata vừa cuống quýt né tránh vừa nhìn hai cô gái vừa đến, gương mặt khổ sở. Cô định nói nhưng chỉ kịp thở dốc: ``Để em kể lại mọi chuyện cho hai người nghe\dots\ Nếu em còn sống.''

\begin{center}
    \uline{Hồi tưởng\\Cách đây không lâu.}
\end{center}

``Khi mà em đã biến thành 3D thì\dots''  cô nàng bắt đầu kể ngắn gọn sự việc, tay vẫn không ngừng né tránh vầng hào quang đang bay vèo vèo.

Tôi từ đó tiếp lời, thở dài với vẻ mặt đầy ân hận: ``\dots\ một quả bóng mà anh ném bâng quơ trúng vào cái vầng hào quang của em ấy\dots\ Anh chỉ muốn thử xem nó có đàn hồi không, ai ngờ\dots''

\begin{figure}[H]
    \centering
    \begin{subfigure}[t]{0.45\textwidth}
        \centering
        \animategraphics[width=8cm]
        {30}{../../../../Image/Book5/s-8g/s-8g.}{1}{60}
        \caption{Khi quả bóng đụng phải vầng hào quang của Kanata…}
    \end{subfigure}
    \quad
    \begin{subfigure}[t]{0.45\textwidth}
        \centering
        \animategraphics[width=8cm]
        {30}{../../../../Image/Book5/s-8g/s-8g.}{61}{161}
        \caption{\dots\ lập tức nó hóa thành khổng lồ.}
    \end{subfigure}
\end{figure}

``\dots\ và thế là nó hóa to luôn! Em không biết tại sao, nhưng nó cứ lớn mãi, xong rồi quay như cối xay gió.''

\begin{center}
    \uline{Kết thúc hồi tưởng.}
\end{center}

Tôi chỉ vào vầng hào quang của cô nàng vẫn đang xoay tít kia, giọng như một người thuyết trình trong chương trình khoa học viễn tưởng: ``Và thế là nó cứ quay như vậy đấy. Không có dấu hiệu dừng lại. Em ấy đã cố gắng bấm vào nó, hét vào nó, thậm chí cầu xin nó, nhưng không ăn thua.''

Aqua và Suisei nhìn nhau, rồi đồng loạt lắc đầu với vẻ hoài nghi. Aqua khoanh tay, khẽ lắc đầu: ``Nghe có vẻ khó tin thật đấy\dots\ Một quả bóng mà biến cái hào quang thành máy bay? Nó như chuyện cổ tích ấy.''

Suisei thì chống cằm, cố tìm ra một lý do hợp lý hơn, như một nhà khoa học đang cố giải thích một hiện tượng lạ: ``Làm gì có chuyện đó! Chắc là cô ấy vô tình bấm vào nút điều chỉnh kích thước trong phần mềm 3D thôi.''

Sau khi nghe hai cô nàng tỏ vẻ không tin với câu chuyện mà cậu và Kanata nói, tôi cau mày, giọng đầy bực bội: ``Mấy đứa nghĩ anh bốc phét sao? Anh có cần lấy camera quay lại cho xem không?''

Aqua liếc về vầng hào quang hình phi tiêu đang vèo vèo bay qua đầu cô, rồi liếc về tôi, nghiêng mặt: ``Không phải là em không tin đâu, nhưng mà\dots\ Ờ thì, vẫn thấy khó tin. Em có thể thử chạm vào nó được không?''

Trong khi đó, Suisei thì hào hứng hơn, mắt sáng rực như trẻ con, cô chỉ tay về vầng hào quang: ``Chị có thể thử đu lên ngôi sao kia được không? Nếu nó quay nhanh thế, có thể dùng nó làm trò chơi cảm giác mạnh!''

Vừa dứt lời, Kanata gần như bốc hỏa. Cô dừng lại, quay phắt người, vầng hào quang suýt cắt ngang tóc Suisei. Mắt cô sáng rực vì phẫn nộ: ``\textbf{Em đã bảo nó nguy hiểm rồi mà! Hai chị muốn gây sự với em hả!?}''

Tôi lập tức can thiệp trước khi mọi thứ đi quá xa, vì tôi đã thấy đủ các cuộc chiến của Kanata với Coco, và không muốn thêm một cuộc chiến mới: ``Thôi thôi, bình tĩnh\dots\ Tính nóng như kem thế\dots''

Câu nói vừa thốt ra, cả ba cô gái lập tức dừng lại. Họ quay sang nhìn tôi, mắt mở to. Suisei đưa tay lên môi, Aqua nhướng mày, Kanata thì chớp mắt. Cả ba cô gái đồng thanh, giọng đầy tò mò: ``Tính nóng như kem là sao!?''

Nghe vậy, tôi gãi đầu, cố gắng tìm một lời giải thích cho câu đùa mà tôi vừa nói, mà bây giờ tôi mới nhận ra là mọi người không hề hiểu theo nghĩa tiếng Nhật thông thường: ``Ờm\dots\ Thì\dots\ trông có vẻ lạnh nhưng bên trong lại rất dễ chảy\dots\ Ý anh là, ngoài mặt thì em ấy có vẻ bình tĩnh, nhưng thực chất rất dễ nổi nóng. Giống như kem vậy.''

Suisei nhìn tôi với ánh mắt đầy thương hại: ``Anh có thể nghĩ ra cách so sánh nào tệ hơn không?''

Aqua cười khúc khích, còn Kanata thì đỏ mặt, vừa xấu hổ vừa tức giận.

Mặc dù cuộc tranh luận về ``tính nóng như kem'' vẫn chưa ngã ngũ, nhưng vầng hào quang sau lưng Kanata vẫn tiếp tục quay với tốc độ đáng sợ, đe dọa bất cứ ai lỡ dại đứng quá gần. Nó vèo qua đầu Aqua, rồi lại lao về phía Suisei, khiến cả hai phải cúi xuống né tránh.

Đến cuối cùng, bằng một cách nào đó, có thể là do Kanata đã bấm nhầm nút, hoặc do vầng hào quang tự động tắt sau một thời gian, nó quay trở về bình thường. Nhưng đó sẽ là một trong những thứ mà mọi người không nên động vào. Tôi sẽ nhớ mãi bài học: đừng bao giờ ném bóng vào vầng hào quang của thiên sứ, và đừng bao giờ so sánh tính cách với kem.

Tôi thở phào nhẹ nhõm, ngồi xuống ghế, nhìn Kanata đang thu dọn hào quang: ``Lần sau nếu có ai hỏi anh về vầng hào quang của em, anh sẽ nói là nó có tính cách riêng.''

Kanata quay sang, nhướng mày: ``Tính cách gì cơ?''

``Tính cách\dots\ bảo vệ chủ nhân khỏi bóng bay.''

Cả phòng im lặng, rồi bật cười.

\subsection*{\phantomsection\label{ep-S7.IV}}
\addcontentsline{toc}{subsection}{{\sen{-Histoire 4-}} {\caviar Ma Sói}}
\stpt{-Histoire 4-}
\stcap{Ma Sói}
\vspace{-4pt}
\begin{center}
    {\caviar\fontsize{12pt}{12pt}\selectfont Nhân vật: \emp{kuon1} \emp{miko2} \emp{pekora} \emp{korone} \emp{marine} \emp{kanata}}
\end{center}

Một buổi chiều định mệnh, khi mà những trò chơi board game tưởng chừng vô hại lại biến thành chiến trường của những nghi ngờ và phản bội. Bảy con người - Kanata, Pekora, Miko, Korone, Marine, và tôi, cùng với Fukuyasu trong vai quản trò - đang ngồi quanh bàn, căng thẳng nhìn nhau trong ván Ma Sói định mệnh.

Trò chơi đã diễn ra được một lúc. Những lời buộc tội đã bay qua lại như tên lửa. Nhưng đến phiên điều trần của mọi người, không ai bảo ai, tất cả đều đồng loạt chỉ tay về phía một người: nàng Thiên sứ tóc xám.

Kanata ngồi đó, mặt hơi xanh, mồ hôi lấm tấm trên trán.

Nàng Thuyền trưởng là người đầu tiên lên tiếng, giọng như một thẩm phán tuyên án, mắt nheo lại đầy nghi ngờ: ``Kanatan, em là Ma Sói sao? Nãy giờ em im như thóc, chẳng nói chẳng rằng, trong khi mọi người tranh luận sôi nổi. Kẻ có tội thường ít nói nhất.''

Nàng Thỏ gật gù, tay chỉ vào Kanata: ``Nãy giờ không thấy nói gì ha, peko! Lúc bình thường em nói như vẹt cơ mà, giờ lại im thin thít. Đúng là đáng nghi.''

Nàng Vu nữ cũng hùa theo, giọng đầy tự tin như thể đã có bằng chứng: ``Chị cũng thấy em là người lén lút nhất đó ne. Chị đã theo dõi em từ đầu ván, em lúc nào cũng nhìn xuống bàn, tránh mắt người khác. Khác hẳn thường ngày.''

\img{s-8h}

Thấy cả ba cô gái đều chỉ điểm vào cô, Kanata tái mét. Cô nhìn quanh, từ người này sang người khác, rồi cố gắng bào chữa, giọng run run: ``Ơ? Nhưng em là dân làng mà\dots\ Em không phải Ma Sói. Em chỉ đang cố gắng suy nghĩ xem ai mới là kẻ địch thôi\dots''

Tôi ngồi ngay bên cạnh, khoanh tay nhìn cô với vẻ mặt đầy nghi vấn, giọng chắc nịch: ``Dân làng mà sao nãy giờ câm như hến thế? Lại còn run như cầy sấy kia kìa? Em biết không, trong Ma Sói, dân làng thường hay nói nhiều, còn Ma Sói thì lại im lặng để tránh lỡ lời. Em đang làm đúng chiến thuật đó.''

Nàng Khuyển ngồi đối diện ngả người ra sau ghế, nở một nụ cười đầy bí ẩn. Cô nhìn Kanata, rồi phán một câu như thể đã đoán được danh tính thật của cô từ lúc bắt đầu trò chơi: ``Còn lời trăn trối cuối cùng nào không, Kanata-chan? Đã đến lúc phải đối mặt với số phận rồi. Linh cảm của loài Chó mách bảo chị rằng em chính là Ma Sói.''

Nàng Thiên sứ chỉ có thể nuốt khan, mắt đảo liên tục giữa những ánh nhìn đầy nghi ngờ. Cô mở miệng định nói gì đó, rồi lại khép lại. Cô nhìn tôi, rồi nhìn Pekora, rồi nhìn Korone, rồi nhìn Marine. Ánh mắt của họ không ai cho cô một cơ hội nào. Đáp lại là một sự im lặng đáng ngờ. Một sự im lặng mà chỉ có người bị buộc tội mới hiểu.

Rồi cả bốn người còn lại phá lên cười giòn giã. Tôi lắc đầu thở dài trong bất lực: ``Đúng là dễ mắc lừa mà. Mới có vài phút nghi vấn đã lộ hết rồi.''

Cười đã đời xong, Marine đập bàn kết luận, giọng như một vị tướng ra lệnh hành quyết: ``Nào, chắc chắn ả ta là Ma Sói! Cùng giết ả ta đi! Số phiếu thuận: bốn. Không có ai bỏ phiếu trắng!''

Pekora vỗ tay bôm bốp tỏ ý đồng tình, mắt sáng lên như thể vừa tìm ra được kẻ thù: ``Thật sự trông em đáng nghi thật đó, peko! Lần này em hết đường chạy rồi.''

Kanata liếc về phía tôi, xem tôi như là một cứu cánh cuối cùng trong cơn bão tố. Mắt cô long lanh, như thể đang cầu xin sự thương xót. Nhưng đáp lại, tôi chỉ nhún vai, giọng đầy vẻ ``xin lỗi'': ``Rất tiếc là\dots\ anh phải đồng ý với mọi người ở đây. Em đã tự đào hố cho mình rồi.''

Miko giơ tay lên trời, giọng vui sướng như thể vừa chiến thắng: ``Vậy là ta chốt ne! Ma Sói bị bắt!''

Không chờ Kanata kịp phản biện, tất cả đã biểu quyết loại cô khỏi vòng chơi. Bốn lá phiếu, một màu đỏ, tất cả đều chỉ về cô.

Khi trời tối trong trò chơi, cả dân làng đi ngủ, chỉ còn nàng Thiên sứ ``ra đi'' trong uất ức. Cô bị loại theo quyết định của Korone, người đã đưa ra lý do duy nhất là\dots\ ``linh cảm của loài Chó''.

\gif{s-8i}{1}{134}

Kanata ngồi ngoài cuộc, mặt đầy phẫn uất, tay nắm chặt, nhìn những người còn lại tiếp tục chơi.

Sáng hôm sau, quản trò Fukuyasu gọi mọi người dậy, giọng đều đều như thường lệ: ``Mọi người ơi, trời sáng rồi, dậy thôi! Hôm qua chúng ta có một người mới chết nữa, đó là Kanata-san.''

Cậu bạn chậm rãi mở tờ giấy trước mặt, rồi thông báo, giọng đầy bất ngờ: ``Và cô ấy\dots\ chính là Ma Sói.''

Cả bốn cô gái lập tức vỡ òa trong sung sướng. Pekora nhảy dựng lên, Marine vỗ tay rầm rầm, Miko ôm mặt cười, Korone thì tự hào vỗ ngực: ``\textbf{Thắng rồi!!!} Linh cảm của tôi không bao giờ sai! Lần đầu tiên trong lịch sử Ma Sói của \hll, dân làng thắng!''

Đây có lẽ là lần hiếm khi mà họ thắng được một ván Ma Sói, bởi vì thường thì ai cũng đoán sai và loại nhầm người.

Trong khi mọi người vui vẻ ăn mừng, Kanata chỉ có thể trừng mắt, lườm nguýt cả bọn đang reo hò như thể vừa giành cúp vô địch World Cup. Cô nàng nổi gân xanh, nước mắt lưng tròng vì cay cú: ``Lần sau đừng có bảo là tôi dễ bị lừa nữa đấy!!!''

Tôi chỉ cười, khoanh tay nhìn cô: ``Anh không nói em dễ bị lừa. Anh nói em không biết giữ bí mật. Đó là hai chuyện khác nhau đấy.''

Kanata chỉ biết đập tay xuống bàn, vừa tức vừa buồn cười, còn mấy người kia thì vẫn tiếp tục ăn mừng.

\subsection*{\phantomsection\label{ep-S7.V}}
\addcontentsline{toc}{subsection}{{\sen{-Histoire 5-}} {\caviar Phiếu bầu}}
\stpt{-Histoire 5-}
\stcap{Phiếu bầu}
\vspace{-4pt}
\begin{center}
    {\caviar\fontsize{12pt}{12pt}\selectfont Nhân vật: \emp{kuon1} \emp{kanata}}
\end{center}

Tôi hắng giọng, nhìn vào tờ giấy kết quả từ cuộc bình chọn mà\dots\ chính Kanata lập ra. Tôi vẫn chưa hiểu lắm vì sao lại có cuộc bình chọn này, nhưng với tư cách là người dẫn chương trình bất đắc dĩ, tôi quyết định không hỏi thêm cho đỡ đau đầu. Có những bí mật mà thiên sứ giữ kín, và tôi không muốn là người khám phá ra chúng.

``Và bây giờ là kết quả cuộc bình chọn nhân vật yêu thích nhất lần đầu tiên từ buổi debut 3D của Amane Kanata!''  tôi chìa tay về phía màn hình, cố gắng dẫn dắt với sự chuyên nghiệp nhất có thể, dù trong lòng tôi đang tự hỏi tại sao mình lại đồng ý làm việc này.

Trên màn hình, những hiệu ứng lấp lánh màu xanh nước biển xuất hiện, cùng với hai dòng chữ:

\begin{center}
    ``hololive production''
\end{center}
và
\begin{center}
    ``Buổi ra mắt 3D đặc biệt của Amane Kanata!!
    Cuộc bình chọn nhân vật yêu thích nhất lần đầu tiên''
\end{center}

Sau vài giây hiệu ứng lòe loẹt (đến mức tôi suýt bị chói mắt), màn hình bắt đầu nhiễu nhẹ rồi hiển thị danh sách năm người đứng đầu. Mỗi người được xếp hạng kèm theo tên và số lượt bình chọn. Tôi nhìn xuống tờ giấy, rồi nhìn lên màn hình, rồi lại nhìn xuống tờ giấy.

``Bắt đầu với\dots\ Amane Kanata ở vị trí thứ 5, với 500 lượt bình chọn!''

Trên màn hình, Kanata viền xanh lá cây đang tạo dáng chỉnh lại cặp kính, rồi thở dài với vẻ mặt hơi buồn: ``Ừm\dots\ kết quả này cũng bình thường thôi\dots\ Ít nhất cũng có người bỏ phiếu cho mình.''

Tôi tiếp tục nhìn vào danh sách, rồi tiếp tục dẫn\dots\ để rồi tôi cảm thấy có gì đó không đúng. Tôi nheo mắt, nhìn kỹ dòng chữ thứ tư: ``Kế đến là\dots\ Amane Kanata? Với 750 lượt!?''

Kanata viền đỏ - một phiên bản khác của cô - cảm thấy bất mãn với kết quả này. Cô nàng tức giận thốt lên, đập tay xuống bàn: ``Híc\dots\ mình thua Kanata mất rồi\dots\ Sao lại có người thích Kanata hơn mình chứ?''

Tôi ngơ ngác nhìn danh sách một lần nữa, để rồi ngạc nhiên với cái tên tiếp theo ở vị trí thứ ba. Tôi đưa tay dụi mắt, nhưng con số vẫn không thay đổi: ``Tiếp theo là\dots\ ủa, Amane Kanata nữa!? Với 950 phiếu\dots?''

Kanata viền tím thở dài chán nản, mắt lơ đễnh thất vọng, nhưng rồi chắp tay lại như cầu nguyện: ``Tạ ơn Chúa\dots\ ít nhất cũng còn được hạng ba. Nhưng mà 950 phiếu so với 750 thì cũng chẳng hơn gì mấy.''

Tôi bắt đầu nhíu mày, cảm giác có gì đó không ổn. Tôi lật tờ giấy kiểm tra lại, nhìn kỹ từng dòng, để rồi ngỡ ngàng khi liếc sang người tiếp theo ở vị trí thứ hai: ``Ờm, kế đến sẽ là\dots\ thật luôn? Amane Kanata, với 952 phiếu\dots''

Trên màn ảnh, nàng Thiên sứ viền xanh đậm khoanh tay, ngoảnh mặt đi, chỉ ``hừm'' một tiếng tỏ vẻ kiêu ngạo: ``Chỉ hơn có 2 phiếu so với hạng ba. Mà thôi, vẫn là thắng.''

Tôi tỏ vẻ bất lực với danh sách này. Năm người đứng đầu, tất cả đều là Kanata. Cùng một tên, cùng một khuôn mặt, chỉ khác màu viền. Tôi cố gắng hoàn thành trách nhiệm với tư cách là người dẫn chương trình, hít sâu, rồi tiếp tục với vị trí số một: ``Và giải nhất sẽ là\dots''

Lập tức, tôi sụp tinh thần hoàn toàn. Tôi ném tờ giấy lên trời, thốt lên bất lực: ``\textbf{Thôi nghỉ cmn đi! Mọi người cũng đoán ra rồi đấy!}''

Tôi đứng dậy, hậm hực bỏ vào cánh gà, mặc cho màn hình tiếp tục chạy.

Trên màn hình, một Kanata viền xanh lam vui vẻ hiện ra với \(10^6-1=999.999\) phiếu, tới mức tràn cả ra ngoài ô. Con số này nhiều đến nỗi nó vượt cả khung hình, như thể đang cố gắng thoát ra ngoài màn hình.

Nàng Thiên sứ phiên bản chiến thắng cười rạng rỡ, vẫy tay đầy tự hào, giọng như một người vừa đoạt giải thưởng lớn: ``Cảm ơn mọi người rất nhiều! Cảm ơn tất cả các Kanata đã bỏ phiếu cho tôi!''

Cảnh tượng còn chưa kết thúc, một dòng chữ to tướng hiện lên trên màn hình, như thể đang chế giễu người xem: ``\textbf{Nhân vật bạn yêu thích nhất đứng ở vị trí nào?}''

\img{s-8j}

Từ trong cánh gà, giọng tôi vang lên đầy chán nản, như một linh hồn đã từ bỏ cuộc sống: ``Tuyệt lắm, lại còn có câu hỏi này nữa chứ! Game này em chơi một mình luôn đi! Cần gì phải có người dẫn chương trình khi mà người chiến thắng đã rõ từ đầu cơ chứ!''

Ngay lập tức, hiệu ứng nổ xuất hiện trên màn hình - những bông hoa, những ngôi sao, và đủ thứ lấp lánh - rồi màn hình chạy chữ kết thúc đoạn phim. Một dòng chữ nhỏ xuất hiện ở góc dưới: ``Kết quả cuộc bình chọn không có giá trị pháp lý. Mọi khiếu nại xin gửi về Kanata.''.

Tôi ngồi trong cánh gà, nhìn ra ngoài, thấy Kanata đang vui vẻ nhảy nhót với chính mình trên màn hình: ``\textit{Đúng là một cuộc bình chọn công bằng\dots\ Nếu công bằng có nghĩa là tất cả phiếu bầu đều là cho cùng một người.}''

\subsection*{\phantomsection\label{ep-S7.VI}}
\addcontentsline{toc}{subsection}{{\sen{-Histoire 6-}} {\caviar Giới thiệu tập sau}}
\stpt{-Histoire 6-}
\stcap{Giới thiệu tập sau}
\vspace{-4pt}
\begin{center}
    {\caviar\fontsize{12pt}{12pt}\selectfont Nhân vật: \emp{kuon1} \emp{haato} \emp{coco} \emp{kanata} \emp{suisei} \emp{miko2} \emp{aqua} \emp{pekora} \emp{korone} \emp{marine}}
\end{center}

\pov{-}

Màn hình tối sầm. Rồi một ngọn lửa bùng lên, chiếm trọn màn ảnh, đỏ rực và mãnh liệt như thể đang thiêu đốt cả thế giới.

Từ phía sau, giọng Kanata vang lên, đầy kịch tính, như một nhân vật chính trong phim hành động sắp đối mặt với số phận: ``Dừng lại đi! Con quái vật hướng ngoại Subaru-senpai bảo tôi rằng\dots''

Giọng cô bất ngờ chuyển sang trầm khàn, bắt chước giọng Vịt của Subaru: ``\dots\ ngày mai chúng ta sẽ ăn gì đi!''

Cảnh quay lần lượt chuyển nhanh qua cận mặt của Aqua, Suisei, Marine, Pekora, Korone, Miko, Coco, Haato và cuối cùng là Kuon, tất cả đều mang vẻ mặt đầy nghiêm trọng: mắt nheo lại, môi mím chặt, như thể đang chuẩn bị cho một trận chiến sinh tử.

Nhạc nền gay cấn dồn dập, tiếng trống đập mạnh như nhịp tim của một người sắp bước vào cuộc chiến. Kanata tiếp tục với chất giọng trầm hẳn xuống, như thể đang tuyên bố số phận của chính mình: ``Nhưng kẻ hướng nội Amane Kanata đây sẽ phải chuẩn bị tinh thần để ra ngoài với bọn họ\dots''

\img{s-8l}

Màn hình chuyển cảnh. Một bóng dáng quen thuộc xuất hiện giữa biển lửa: cậu Tổng, tay chống hông, nhìn vào cô nàng đang run rẩy trong ngọn lửa, lẩm bẩm với giọng đầy bất lực: ``Cô ta đã hóa thành tro rồi\dots''

Kanata tiếp tục tự độc thoại, giọng đầy bi tráng, như một chiến binh sắp hy sinh vì chính nghĩa: ``Xin đừng chết, Kanata à! Nếu cậu chết ở đây\dots\ Lời hứa đi ăn mì ramen cùng với Coco sẽ ra sao đây!?''

Cảnh cận mặt Kanata, đôi mắt cô ánh lên một tia quyết tâm mãnh liệt. Nhạc nền dồn dập hơn, như thể đang đẩy cao cảm xúc: ``Mình vẫn còn chút thời gian\dots\ Nếu mình có thể sống sót được buổi giới thiệu đầu tiên, mình có thể sống sót được những lần tiếp theo\dots''

\img{s-8m}

Màn hình tối sầm lại, rồi bật sáng với dòng chữ hoành tráng, sắc đỏ như máu:

\begin{center}
    \textbf{🔥 CÁI CHẾT CỦA KANATA 🔥}
\end{center}

Ngay bên cạnh tiêu đề là gương mặt cười khẩy của nàng tóc trắng, đầy bí hiểm như thể nắm giữ lá bài lật ngửa trong tay. Trông cô vừa đáng sợ, vừa buồn cười, giống như một nhân vật phản diện đang cố gắng tỏ ra nguy hiểm nhưng thất bại thảm hại.

Và rồi, một giọng nói trầm hùng vang lên, vọng khắp không gian: ``\textbf{Dual Standby!}''

Màn hình vụt tắt. Không gian bỗng lặng như tờ.

Im lặng kéo dài vài giây. Rồi cậu Tổng - người đã đứng ngoài cuộc từ đầu - cuối cùng cũng lên tiếng, giọng đầy bất lực, phá tan bầu không khí kịch tính: ``\dots Anh nghĩ là chúng ta làm hơi quá rồi đấy. Đi ăn thôi mà sao như đi đánh trận vậy\dots\ Không biết ai đã viết kịch bản cho cái trailer này, nhưng chắc chắn người đó đã xem quá nhiều phim hành động.''

Cảnh quay kết thúc bằng hình ảnh chữ ``Dual Standby'' mờ dần, như một lời hứa hẹn cho những điều sắp tới.

\subsection*{\phantomsection\label{ep-S7.VII}}
\addcontentsline{toc}{subsection}{{\sen{-Histoire 7-}} {\caviar Bữa ăn}}
\stpt{-Histoire 7-}
\stcap{Bữa ăn}
\vspace{-4pt}
\begin{center}
    {\caviar\fontsize{12pt}{12pt}\selectfont Nhân vật: \emp{kuon1} \emp{haato} \emp{coco} \emp{kanata} \emp{suisei} \emp{aqua} \emp{miko2} \emp{pekora} \emp{korone} \emp{marine}}
\end{center}

\pov{Hikawa Kuon}

Buổi tối tại quán mì ramen quen thuộc, ở cái quán nhỏ nằm nép mình trong một con hẻm mà chỉ những người trong công ty mới biết. Chín con người quây quần bên bàn gỗ, hơi nóng bốc lên từ những tô mì, mang theo hương thơm đậm đà của nước dùng xương hầm. Tiếng đũa gõ vào bát, tiếng xì xụp mì, tiếng cười nói rôm rả - tất cả tạo nên một bầu không khí ấm cúng, khác hẳn với cái sự ồn ào ngoài kia.

Phải thuyết phục mãi, Kanata mới chịu rời khỏi cái kén an toàn của mình để tham gia cùng nhóm. Tôi đã phải nhờ Coco làm ``lực lượng đặc biệt'' mới kéo được cô ra khỏi văn phòng. Nhưng vì là Coco mời, và còn vì cô ấy hứa sẽ trả tiền, nàng Thiên sứ quyết định vứt bỏ mọi lo âu mà ăn một cách xả láng, mặc kệ số phận.

Đũa của cô nàng Thiên sứ lia lịa, cô gắp từng miếng thịt, từng sợi mì, rồi húp nước súp một cách đầy hài lòng. Trông cô như một chiến binh vừa chiến thắng một trận đánh lớn.

Cả nhóm đang tập trung tận hưởng hương vị của món ăn, thỉnh thoảng lại có tiếng ``ngon quá'' của Suisei, tiếng ``peko'' của Pekora, tiếng cười của Korone - thì tôi đặt bát mì xuống, nhướng mày nhìn Kanata: ``Cái Y-G-O mà em vừa làm ban nãy có cần thiết không? Trông như trailer của một bộ phim thảm họa ấy. Ai mà biết đi ăn mì lại cần có kịch bản hành động?''

Kanata đang xì xụp bát mì bỗng khựng lại, rồi ngẩng đầu lên. Đôi mắt cô ánh lên sự tự hào, như một đạo diễn vừa hoàn thành tác phẩm của mình: ``Em thấy ngầu nên cứ làm thôi! Mà cũng có tác dụng mà, cuối cùng cũng kéo được mọi người ra ngoài.''

Suisei bật cười, vừa gắp mì vừa nói: ``Ngầu thì cũng được đấy, nhưng mà chị thấy giống một đoạn trailer anime hơn! Phim hành động của Nhật Bản, có lửa, có thoại kịch tính, có cận mặt, chỉ thiếu mỗi nhạc nền của một ban nhạc rock thôi.''

Pekora vừa nhai sợi mì vừa gật gù, tay cầm đũa chỉ về phía Kanata: ``Phải đó, cứ như đang quảng cáo cho một bộ shonen nào đó ấy, peko! Mà nếu có thêm phần 'mùa sau' thì càng hay.''

Aqua ngồi bên cạnh, vừa thổi nguội tô mì vừa góp lời, giọng đầy vẻ suy tư của một chuyên gia ẩm thực: ``Chị thấy nó giống trailer của một game hành động hơn. Kiểu như mấy game mà nhân vật chính phải chiến đấu để giành lại bữa tối ấy!''

Marine chống cằm, cười xảo quyệt, giọng đầy tinh nghịch: ``Chị cá là ngày hologra ra mắt, tụi fan sẽ spam 'DRAW' với 'STANDBY' khắp nơi cho xem. Có khi nó thành trend mất.''

Cả bàn rôm rả cười đùa, không khí vô cùng sôi động. Ngay cả Korone, người vốn chỉ tập trung vào ăn, cũng phải ngẩng lên cười theo.

Miko bất ngờ quay sang hỏi Coco, vẻ mặt đầy tinh nghịch: ``Bữa này em đãi hả, ne? Chị thấy em gọi món đắt nhất trên menu đấy.''

Coco nhếch môi, vỗ ngực đầy tự tin, giọng như một đại gia sắp chi trả hóa đơn: ``Đương nhiên! Chị phải cổ vũ Kanatan chúng ta chứ! Dù sao cũng là lần đầu em ấy chịu ra ngoài ăn với mọi người.''

Kanata nhìn sang Coco, cảm động nhưng cũng có chút lo lắng, vì cô biết rằng Coco có thể vung tiền nhưng cũng có thể quên ví ở nhà: ``Vậy có cần tôi trả lại sau không?''

Coco cười phá lên, vỗ mạnh vào lưng Kanata khiến cô suýt sặc nước súp. Cô nàng Rồng nói, giọng đầy khích lệ: ``Không cần! Nhưng nhớ là ngày hôm đó phải làm tốt đấy! Đừng để tôi đây thất vọng, hiểu chưa? Làm tốt thì bữa sau chị lại đãi tiếp.''

Kanata gật đầu thật mạnh, ánh mắt lấp lánh đầy quyết tâm, như một chiến binh sẵn sàng ra trận. Có lẽ cô đang nghĩ đến việc chuẩn bị cho buổi hologra, hoặc có lẽ chỉ đơn giản là nghĩ đến món mì tiếp theo.

Tám người con gái còn lại cũng lần lượt cổ vũ cô nàng. Họ gửi lời chúc cho Kanata có một buổi ra mắt hologra thật tốt, nơi cô sẽ thực sự tỏa sáng\dots\ cùng với đầy sự rắc rối và hỗn loạn.

Suisei giơ cốc nước lên: ``Chúc Kanata-chan có một buổi ra mắt thật thành công!''

Aqua tiếp lời, giọng đầy nhiệt tình: ``Và đừng quên mỉm cười với camera thật tươi nhé! Chị nghe nói người ta thích idol cười tươi nhất!''

Haato vẫy tay: ``Và đừng quên mỉm cười thật tươi nhé!''

Marine cười đểu: ``Và đừng có ngã trên sân khấu đấy!''

Cả bàn cùng cười, tiếng đũa, tiếng bát, tiếng nói chuyện hòa lẫn thành một bản giao hưởng của niềm vui.

Và theo đúng lịch trình, đó cũng sẽ là ngày nàng Thiên sứ gặp lại Coco và tôi trên sân khấu, cùng với lần đầu tiên chạm mặt nàng giáo viên quyến rũ Choco-sensei\dots\ và nàng Quỷ sừng tinh quái với câu nói thương hiệu: ``Yo-dayo!''

Tôi nhìn quanh bàn, nhìn những khuôn mặt đang cười đùa, rồi nhìn Kanata đang cúi xuống tô mì, và thở dài: ``\textit{Ít nhất thì lần này cũng không có bom nổ. Không có vẹt. Không có hào quang to như phi tiêu. Chỉ là một bữa ăn bình thường. Đúng là một ngày thành công.}''

Rồi tôi cũng cầm đũa lên, hòa vào không khí ấm áp của buổi tối.

\begin{center}
    \uline{(Với những bạn nào đã đọc đến đây, quay lại {\flaregothic ÉPISODE 62} để đọc (lại) câu chuyện trên.)}
\end{center}

\end{document}